% === §7 Plan de trabajo ===
% Redactado por el Redactor — Fase 5
% Alineación estricta con las fases de la Metodología (§6): F1→F2→F3→F4
% Duración: 18 meses de ejecución + 2 meses de productos = 20 meses

% --- Párrafo introductorio ---

El plan de trabajo se estructura en correspondencia directa con las cuatro fases
de la metodología (sección~\ref{sec:metodologia}): F1 (fundamentación y
caracterización del aprendiz, OE1), F2 (diseño y desarrollo del ecosistema
agéntico, OE2), F3 (validación empírica en aula y TRL~6, OE3) y F4 (transferencia
y proyección TRL~7, OE3). La duración total se ajusta a los términos de la
Convocatoria Nacional Alianzas SIUN 2025--2027: \textbf{18 meses de ejecución más
2 meses adicionales para productos, totalizando 20 meses}. Las fases presentan
solapes intencionales que reflejan el carácter iterativo e incremental del
desarrollo (cadencia tipo Scrum híbrido): F1 ocupa los meses 1--5, F2 los meses
4--12, F3 los meses 11--16 y F4 los meses 15--20, de modo que las fases
consecutivas transfieren entregables sin interrupciones y el riesgo de cuello de
botella se reduce. Cada actividad designa un rol responsable y culmina en un hito
verificable en el que se libera un producto tangible, conectando con los
resultados de la sección~\ref{sec:resultados}.

\medskip

% --- Descripción del cronograma Gantt para el Diseñador ---
% NOTA PARA EL DISEÑADOR (agente TikZ): implementar con pgfgantt en
% proposal/sections/diag_gantt.tex.

\noindent\textbf{Descripción del diagrama de Gantt.}
El cronograma se renderizará como un \emph{chart} de \texttt{pgfgantt} con el
eje horizontal graduado en 20 meses (segmentado visualmente: meses 1--18 de
ejecución sobre fondo neutro, meses 19--20 de productos con un tono diferenciado).
Cada fase se representa como una barra horizontal que abarca sus meses de inicio y
fin: F1 (meses 1--5), F2 (meses 4--12), F3 (meses 11--16) y F4 (meses 15--20).
Las barras deben usar la paleta institucional (azul UNAL, gris LabIA, verde
GCPDS) asignando un color por fase. Sobre cada barra se etiquetan las actividades
principales como sub-barras \emph{linked} (encadenadas) con su intervalo mensual
específico (ver tabla siguiente). Se deben marcar cuatro \textbf{hitos} como
nodos simbólicos sobre el eje temporal: H1 (mes 5, TRL~4), H2 (mes 12, TRL~5), H3
(mes 16, TRL~6) y H4 (mes 20, TRL~7 proyectado), cada uno con una etiqueta y un
icono distintivo (triángulo o diamante). Una cinta paralela inferior debe
representar la progresión TRL~4$\to$5$\to$6$\to$7 alineada con los hitos. El
Diseñador debe asegurar que el gráfico sea legible en escala de grises e incluir
una leyenda con los responsables (IP, coinvestigadores, \emph{team lead},
estudiantes de posgrado, estudiantes de pregrado).

\medskip

% --- Tabla de actividades ---

\begin{table}[h]
\centering
\caption{Plan de trabajo: actividades por fase, responsables, meses y
productos/hitos. Total: 20 meses (18 de ejecución + 2 de productos).}
\label{tab:plan_trabajo}
\small
\begin{tabularx}{\textwidth}{@{}l X l c c X@{}}
\toprule
\textbf{Fase} & \textbf{Actividad} & \textbf{Responsable} & \textbf{Ini.} & \textbf{Fin} & \textbf{Producto / Hito} \\
\midrule

% --- Fase 1 (OE1) ---
F1 & A1.1 Diseño y validación psicométrica del instrumento de diagnóstico de aprendizaje profundo
   & IP + D.\,Angulo + est.\ doctorado & 1 & 3 & Instrumento validado (IVC, $\alpha$ Cronbach) \\

F1 & A1.2 Recolección y curaduría del dataset multimodal de trazas (textos, código, simulación)
   & est.\ doctorado + est.\ pregrado & 2 & 5 & Dataset curado con consentimiento ético \\

F1 & A1.3 Desarrollo y entrenamiento del modelo computacional del aprendiz (E1)
   & IP + est.\ doctorado & 3 & 5 & Modelo del aprendiz (F1, AUC-ROC, F1) \\

\addlinespace
\multicolumn{6}{c}{\textit{Hito H1 (mes 5): TRL~4 — componentes de diagnóstico validados en laboratorio.}} \\
\addlinespace

% --- Fase 2 (OE2) ---
F2 & A2.1 Diseño de la ontología de roles agénticos y arquitectura multiagente (E2)
   & IP + C.\,Castellanos + \emph{team lead} & 4 & 6 & Documento de arquitectura agéntica \\

F2 & A2.2 Implementación del bus de orquestación MCP y memoria compartida
   & \emph{team lead} + est.\ maestría & 5 & 9 & Núcleo de orquestación MCP funcional \\

F2 & A2.3 Desarrollo del módulo de scaffolding adaptativo (E3)
   & est.\ maestría + est.\ pregrado & 6 & 10 & Función $s(\mathbf{x})$ calibrada \\

F2 & A2.4 Implementación del generador de problemas y evaluador formativo (E4)
   & \emph{team lead} + est.\ pregrado & 7 & 10 & Tubería RAG + evaluador NLP \\

F2 & A2.5 Integración del ecosistema y pruebas en entorno simulado
   & IP + equipo completo & 10 & 12 & Prototipo integrado + métricas de desempeño \\

\addlinespace
\multicolumn{6}{c}{\textit{Hito H2 (mes 12): TRL~5 — ecosistema integrado validado en entorno simulado.}} \\
\addlinespace

% --- Fase 3 (OE3, E5) ---
F3 & A3.1 Diseño del protocolo cuasiexperimental y aprobación del Comité de Ética (E5)
   & IP + D.\,Angulo + est.\ doctorado & 11 & 12 & Protocolo aprobado + asignaturas seleccionadas \\

F3 & A3.2 Ejecución de la intervención en aula (grupos experimental y control)
   & est.\ doctorado + 2 est.\ pregrado & 12 & 15 & Datos longitudinales T0--T2 \\

F3 & A3.3 Análisis estadístico pre--post y analítica del aprendizaje
   & est.\ doctorado + IP & 14 & 16 & Informe de impacto causal ($d$ de Cohen) \\

\addlinespace
\multicolumn{6}{c}{\textit{Hito H3 (mes 16): TRL~6 — evidencia causal de impacto en aula real.}} \\
\addlinespace

% --- Fase 4 (OE3, E6) — incluye 2 meses de productos ---
F4 & A4.1 Despliegue en infraestructura híbrida (Apple Silicon + VPS) (E6)
   & \emph{team lead} + est.\ maestría & 15 & 17 & Ecosistema desplegado en prod. \\

F4 & A4.2 Documentación de interoperabilidad, registro de software y código abierto
   & IP + profesional transferencia & 16 & 19 & Registro de software + repo GitHub \\

F4 & A4.3 Plan de adopción con actor del sector productivo (demostración TRL~7)
   & IP + profesional transferencia & 17 & 19 & Acta de demostración operativa \\

F4 & A4.4 Redacción y sometimiento de artículos indexados + productos de apropiación
   & IP + coinvestigadores & 18 & 20 & Artículos Q1/Q2 sometidos + ponencias \\

\addlinespace
\multicolumn{6}{c}{\textit{Hito H4 (mes 20): TRL~7 proyectado — productos académicos entregados.}} \\
\bottomrule
\end{tabularx}
\end{table}

\medskip

% --- Párrafo de hitos y productos ---

\noindent\textbf{Hitos y productos.}
El plan define cuatro hitos verificables al cierre de cada fase, cada uno
asociado a la liberación de productos tangibles que constituyen los resultados de
la sección~\ref{sec:resultados}. \textbf{H1 (mes 5, TRL~4):} entrega del
instrumento de diagnóstico validado psicométricamente, el \emph{dataset}
multimodal curado y el modelo computacional del aprendiz; constituye el primer
producto de generación de nuevo conocimiento y habilita F2. \textbf{H2 (mes 12,
TRL~5):} liberación del prototipo integrado del ecosistema agéntico con métricas
de desempeño documentadas; constituye el primer producto tecnológico
intermedio. \textbf{H3 (mes 16, TRL~6):} entrega del informe de impacto causal
cuasiexperimental con evidencia cuantitativa y cualitativa del efecto del
ecosistema en el aprendizaje profundo; constituye el principal producto de
generación de nuevo conocimiento (artículo Q1/Q2 en preparación) y certifica la
meta de madurez TRL~6 exigida por la convocatoria. \textbf{H4 (mes 20, TRL~7
proyectado):} entrega final durante los dos meses adicionales de productos
—ecosistema desplegado y documentado, registro de software, código abierto,
acta de demostración operativa con un actor del sector productivo, artículos
indexados sometidos y ponencias—; certifica la transferencia tecnológica y la
proyección TRL~7. El solape entre fases y la asignación de responsables por
actividad garantizan que cada producto se libere en el mes previsto sin
depender de una única persona, mitigando riesgos de cuello de botella.